\title{Screening for Bid Rigging with Frequent Losers\\in Public Procurement\thanks{This version: April 2026.}}

\author{Darcio Genicolo-Martins\thanks{Corresponding author. INSPER, S\~ao Paulo, Brazil. Email: darciogm1@insper.edu.br.} \and Paulo Furquim de Azevedo\thanks{INSPER, S\~ao Paulo, Brazil.}}

\date{}

\maketitle

\begin{abstract}
\noindent
Enforcement agencies investigating bid rigging need triage tools
that work with limited data; most existing screens require granular
bid amounts unavailable in many procurement systems. We propose a
participation-based screen---\emph{frequent losers} (FL), firms
that never win yet bid abnormally often---requiring only win/loss
records. Applied to 4.5 million tender-items in Brazilian
procurement (2009--2019), the screen achieves AUC~$= 0.94$ against
competition-authority convictions and is nearly orthogonal to
bid-level tools (correlation 0.06). FL-flagged tenders show
3.6--7.7\% higher conditional prices across four estimation
approaches; the association is not causal. Five supporting
diagnostics---including Bajari--Ye bid-coordination tests, network
heterogeneity, and structural estimation---are jointly coherent with
coordinated cover bidding but do not uniquely identify the
mechanism. The screen is a first-stage triage device that flags
suspicious procurement environments for investigation, not a tool
for establishing firm-level liability.
\end{abstract}

\noindent\textbf{Keywords:} bid rigging, cover bidding, public procurement, cartel screening, frequent losers

\noindent\textbf{JEL:} D44, D73, H57, K21, L41

\bigskip
